% ---------------------------------------------------------------------------
\section{Business Understanding and the KPI Framework}
\label{sec:business}
% ---------------------------------------------------------------------------

\subsection{The valuation problem}

A listing platform intermediates between sellers who want a defensible asking price
and buyers who want to know whether that price is reasonable. Neither party
commissions an appraisal at the point of listing, so the platform fills the gap with
a model-based estimate. The estimate is consequential in both directions. If it is
too high, the listing stagnates and the platform absorbs the cost of a failed
transaction and a dissatisfied agent; if it is too low, the seller loses value and
the platform loses the credibility that makes its estimates worth displaying. The
operational question is therefore not merely ``what is the price'' but ``what is the
price, and how much should anyone trust this number''.

Formally, for a property with observable characteristics $\mathbf{x}$ at listing
time, the product must return a point estimate $\hat{y}(\mathbf{x})$ and an interval
$[\ell_\alpha(\mathbf{x}), u_\alpha(\mathbf{x})]$ whose empirical coverage is close to
its nominal level $1-\alpha$. The project treats this as decision support with a
human in the loop: the model frames the price, and a valuer resolves cases in which
the model is not confident.

\subsection{Stakeholders and decisions}

Table~\ref{tab:stakeholders} links each stakeholder to the decision the model
informs and the requirement that decision imposes. Different stakeholders consume
different outputs: the agent needs a band, the seller needs sensitivity to quality and
renovation, the buyer needs to know whether an asking price is defensible, operations
need a publication rule, and the model risk owner needs evidence of continuing
fitness. That plurality is why the KPI framework mixes statistical, operational and
economic indicators rather than relying on a single accuracy metric.

\begin{table}[htbp]
\centering
\caption{Stakeholders, the decisions supported by the valuation product, and the
corresponding modelling requirements.}
\label{tab:stakeholders}
\small
\begin{tabularx}{\linewidth}{@{}L{0.19\linewidth}L{0.30\linewidth}X@{}}
\toprule
\textbf{Stakeholder} & \textbf{Decision supported} & \textbf{Requirement on the model} \\
\midrule
Listing agent & Recommended asking-price band & Point estimate plus a calibrated interval \\
Seller & Whether to list now or renovate first & Sensitivity to quality, age and renovation \\
Buyer & Whether an asking price is defensible & Interval coverage and monitored bias \\
Platform operations & Whether to publish an estimate automatically & Confidence thresholds and escalation policy \\
Model risk owner & Whether the model remains fit for purpose & Drift and performance monitoring with retraining triggers \\
\bottomrule
\end{tabularx}
\end{table}

The workflow follows directly. The agent enters a listing profile; the model returns
a point estimate with both intervals; high-confidence estimates may be published
without review; medium and low confidence estimates are routed to a valuer; and a
monthly monitoring report determines whether the model remains within tolerance. A
sustained breach opens a retraining ticket (Section~\ref{sec:monitoring}).

\subsection{KPIs, targets and achieved values}

Table~\ref{tab:kpis} states each indicator, its definition, the target agreed before
modelling and the achieved out-of-fold value on 1,458 cleaned transactions. Two
design choices require justification. The interval-coverage targets are two-sided,
because an interval that is too wide satisfies coverage while being operationally
useless. The over-valuation rate is targeted at a band rather than at zero, because a
model that systematically predicted below the sale price would score well on a loss
function that penalises over-valuation more heavily while being commercially
unacceptable to sellers.

\begin{table}[htbp]
\centering
\caption{KPI framework with targets fixed before modelling and achieved values from
the selected model. Out-of-fold values are computed from five-fold cross-validation
predictions over 1,458 transactions.}
\label{tab:kpis}
\small
\begin{tabular}{@{}L{0.25\linewidth}L{0.34\linewidth}rr@{}}
\toprule
\textbf{Indicator} & \textbf{Definition} & \textbf{Target} & \textbf{Achieved} \\
\midrule
RMSE & Root mean squared error (\USD) & $\leq$ 25,000 & 20,019 \\
MAE & Mean absolute error (\USD) & $\leq$ 17,000 & 13,537 \\
MAPE & Mean absolute percentage error & $\leq$ 10\% & 7.98\% \\
$R^2$ & Share of price variance explained & $\geq$ 0.90 & 0.9365 \\
Absolute hit rate & Share within $\pm$15,000 \USD & $\geq$ 65\% & 69.7\% \\
Relative hit rate & Share within $\pm$8\% & $\geq$ 60\% & 64.9\% \\
Coverage, 80\% & Sales inside the 80\% interval & 0.75--0.85 & 0.801 \\
Coverage, 90\% & Sales inside the 90\% interval & 0.85--0.95 & 0.901 \\
Over-valuation rate & Share of estimates above the price & 40--60\% & 48.0\% \\
Mean asymmetric cost & Weighted error per listing (\USD) & monitored & 16,708 \\
Median asymmetric cost & Weighted error per listing (\USD) & monitored & 11,358 \\
\bottomrule
\end{tabular}
\end{table}

\subsection{The asymmetric cost model}

An over-priced listing is assumed to cost the platform 1.5 times as much as an
under-priced one, reflecting the observation that a stale, over-priced listing
consumes marketing attention, fails to transact and damages the agent relationship,
whereas a conservative estimate costs the seller a smaller and more recoverable
amount. The loss for one listing is

\begin{equation}
L(\hat{y}, y) \;=\; 1.5 \cdot \max(\hat{y} - y,\, 0) \;+\; 1.0 \cdot \max(y - \hat{y},\, 0),
\label{eq:loss}
\end{equation}

\noindent where $\hat{y}$ is the published estimate and $y$ the realised price. The
out-of-fold mean of Equation~\eqref{eq:loss} is 16,708 \USD{} per listing with a
median of 11,358 \USD. The over-valuation rate of 48.0 per cent is close to balanced,
but the mean signed error is $-853$ \USD, so the model predicts slightly below the
realised price. That direction is favourable under Equation~\eqref{eq:loss}, but it is
a by-product rather than an objective: the model minimises squared error in log space
and the business loss is used only for reporting and for setting review thresholds.

The weights 1.5 and 1.0 are a documented assumption, not an estimated quantity, so
the reported cost is conditional on them. The team verified that the ranking of the
eight candidate models is unchanged at weights of 1.0 and 2.0, so selection does not
depend on the exact asymmetry; the absolute cost figures do. The weights should be
revisited with the platform's finance function before the loss is used in any
commercial commitment.

\subsection{Tolerance, escalation and success criteria}

The absolute tolerance of 15,000 \USD{} is approximately the width of a normal
negotiation margin on the median property of 163,000 \USD, and it defines the
boundary at which an estimate is treated as agent-ready. The relative tolerance of 8
per cent matches the mean absolute percentage error of the selected model, so it
describes delivered accuracy rather than an aspiration. Estimates whose 90 per cent
interval half-width exceeds 30 per cent of the point estimate are classified as low
confidence and routed to a valuer.

The project fixed five success criteria at the outset: a pipeline that reproduces
every artefact from a single command; a model that meets the accuracy and calibration
targets; an API and console whose outputs agree with the offline evaluation; a
monitoring report that detects the simulated drift and recommends retraining; and
documentation of decisions, limitations and business implications. All five are met
except public deployment, which depends on credentials outside the team's control.
The risks identified at the start were target leakage, overfitting on a small
sample, post-deployment shift, segment bias and synthetic-data realism. Each has a
control: an enforced leakage list (Section~\ref{sec:features}), five-fold
cross-validation with regularisation (Section~\ref{sec:models}), drift and coverage
monitoring (Section~\ref{sec:monitoring}), segment-level error reporting and
documented generation logic (Section~\ref{sec:data}).
